% =====================================================================
%  Chapter 3 — Alterations
%  File: chapters/03_alterations.tex
%  Styled to match Chapters 1-2 (boxed environments + Insight callouts).
%  Cross-references Chapter 1 (alteration, 2-colourability baseline) and
%  Appendix A (Markov) instead of repeating them.
% =====================================================================
\providecommand{\E}{\mathbb{E}}
\providecommand{\Prob}{\mathbb{P}}
\providecommand{\ind}{\mathbf{1}}

\chapter{Alterations}
\label{chap:alterations}

The first moment, used bare, asks a random object to be perfect. The
\emph{alteration method} asks only that it be \emph{nearly} perfect, then pays
to repair the few defects. We met it briefly in Chapter~1, where a random
$2$-colouring of $K_n$ was patched into a Ramsey colouring by deleting one
vertex from each monochromatic clique (Theorem~\ref{thm:ramsey_alt}). The
accounting is always the same: bound the expected number of blemishes by the
first moment, subtract it from what you started with, and conclude that some
outcome leaves enough behind. This chapter pushes the idea through four
problems of increasing subtlety --- dominating sets, the Heilbronn triangle
problem, graphs of high girth and high chromatic number, and the best known
bound for two-colourable hypergraphs.

\begin{insight}
The reason alteration so often beats the bare first moment is a division of
labour. A single random construction must simultaneously be large \emph{and}
defect-free, and demanding both at once is wasteful. Alteration decouples them:
make the object as large as you like and tolerate defects, then spend a
controlled amount of that size erasing them. The net size is
(initial size) $-$ (expected repairs), and as long as the repairs are a lower-order
cost, you keep almost everything. Watch for this two-line structure --- build
generously, then delete --- in every proof below.
\end{insight}

% =====================================================================
\section{Dominating sets}
% =====================================================================

\begin{definition}{Dominating set}{dominating}
In a graph $G=(V,E)$, a set $U\subseteq V$ is \emph{dominating} if every vertex
outside $U$ has a neighbour inside $U$.
\end{definition}

\begin{theorem}{Small dominating sets}{dominating-thm}
Every $n$-vertex graph with minimum degree $\delta\ge1$ has a dominating set of
size at most
\[
  \frac{1+\log(\delta+1)}{\delta+1}\;n.
\]
\end{theorem}

\begin{proof}
Fix $p\in[0,1]$, to be chosen. Include each vertex in a random set $X$
independently with probability $p$, and let
\[
  Y = V\setminus\bigl(X\cup N(X)\bigr)
\]
be the vertices that are neither chosen nor adjacent to a chosen vertex. Then
$X\cup Y$ is dominating: any vertex outside it lies in $X\cup N(X)$, hence has a
neighbour in $X$.

A vertex $v$ lands in $Y$ only if $v$ and all of its $\deg(v)\ge\delta$
neighbours avoid $X$ --- that is $1+\deg(v)\ge1+\delta$ independent misses --- so
$\Prob(v\in Y)\le(1-p)^{1+\delta}$. By linearity,
\[
  \E|X\cup Y| = \E|X|+\E|Y|
  \le pn + (1-p)^{1+\delta}n
  \le \bigl(p + e^{-p(1+\delta)}\bigr)n,
\]
using $1-p\le e^{-p}$. The bracket is minimized at
$p=\dfrac{\log(\delta+1)}{\delta+1}$, where $e^{-p(\delta+1)}=\dfrac{1}{\delta+1}$,
giving the value $\dfrac{1+\log(\delta+1)}{\delta+1}$. Some outcome attains at
most the expectation.
\end{proof}

\begin{insight}
The trivial lower bound on a dominating set is $n/(\delta+1)$ --- each chosen
vertex can cover at most $\delta+1$ vertices --- and this theorem says you can
come within a $\log(\delta+1)$ factor of it. A greedy algorithm (repeatedly
take the vertex covering the most as-yet-uncovered vertices) achieves the same
order and is in fact the standard \textsc{Set-Cover} greedy, whose
$\log$ factor is known to be unavoidable in general. What the alteration proof
buys is brevity: the messy, history-dependent greedy analysis is replaced by a
one-line expectation, with the optimal $p$ falling out of a single derivative.
\end{insight}

% =====================================================================
\section{The Heilbronn triangle problem}
% =====================================================================

\begin{question}{Heilbronn's problem}{heilbronn-q}
How should one place $n$ points in the unit square so that the smallest triangle
they span is as large as possible? Writing
\[
  \Delta(n) = \sup_{\substack{S\subseteq[0,1]^2\\|S|=n}}\
  \min_{\substack{p,q,r\in S\\\text{distinct}}} \operatorname{area}(pqr),
\]
how does $\Delta(n)$ behave?
\end{question}

Naive placements do poorly. Points equally spaced on a circle leave a triangle
of area $\Theta(n^{-3})$ (three consecutive points), and a grid does worse
still --- three collinear points give area zero. Heilbronn conjectured
$\Delta(n)=O(n^{-2})$; Koml\'os, Pintz, and Szemer\'edi (1982) disproved it,
showing $\Delta(n)\gtrsim n^{-2}\log n$ by an intricate random construction.
Here is a simpler argument for the weaker --- but already conjecture-refuting in
spirit --- bound $\Delta(n)\gtrsim n^{-2}$.

\begin{theorem}{Many points, no small triangle}{heilbronn-thm}
For every $n$ there is a set of $n$ points in $[0,1]^2$ in which every triple
spans a triangle of area at least $cn^{-2}$, for an absolute constant $c>0$.
\end{theorem}

\begin{proof}
Scatter $2n$ points independently and uniformly in $[0,1]^2$, and estimate the
chance that a fixed triple is nearly collinear. The distance between two random
points satisfies, by an annulus area estimate,
$\Prob(|pq|\in[x,x+dx])\le 2\pi x\,dx$. Conditioning on $p,q$, the triangle
$pqr$ has area at most $\varepsilon$ exactly when $r$ lies within distance
$2\varepsilon/|pq|$ of the line through $p,q$, a strip of area
$\lesssim \varepsilon/|pq|$ inside the square. Hence
\[
  \Prob\bigl(\operatorname{area}(pqr)\le\varepsilon\bigr)
  \;\lesssim\;\int_0^{\sqrt2} 2\pi x\cdot\frac{\varepsilon}{x}\,dx
  \;\asymp\;\varepsilon.
\]
Let $X$ count the triples of area at most $\varepsilon$. There are
$\binom{2n}{3}=O(n^3)$ triples, so $\E[X]=O(\varepsilon n^3)$. Choosing
$\varepsilon=c/n^2$ with $c$ small makes $\E[X]\le n$.

Now \emph{alter}: delete one point from each offending triple. At least
$2n-X$ points survive, with expectation $\E[2n-X]\ge n$, and the survivors
contain no triangle of area below $\varepsilon=c/n^2$. Some outcome therefore
leaves $\ge n$ points with every triangle of area $\ge cn^{-2}$.
\end{proof}

\begin{remark}{An explicit construction}{heilbronn-algebraic}
A clean derandomization, due to Erd\H{o}s, matches $\Delta(n)\gtrsim n^{-2}$
algebraically. For a prime $p$, the parabola points
$\{(x,x^2)\bmod p : x\in\mathbb F_p\}$ contain no three collinear points (a
parabola meets a line in at most two points), so viewed in the integer grid
$[p]^2$, every triangle has area at least $1/2$ by Pick's theorem. Rescaling
into the unit square gives $p$ points with all triangle areas
$\gtrsim p^{-2}$; for general $n$, use a prime between $n$ and $2n$.
\end{remark}

\begin{insight}
The probabilistic proof and the algebraic one make an instructive pair: the
random construction is effortless to set up but yields an object one cannot
write down, while the parabola is fully explicit yet requires the clever
observation that ``no three collinear'' forces area $\ge1/2$ over the integers.
The $\log n$ that Koml\'os--Pintz--Szemer\'edi gained over both comes from a far
more careful random construction in which the deletions are correlated rather
than treated independently --- a recurring theme: the easy alteration leaves a
logarithmic factor on the table.
\end{insight}

% =====================================================================
\section{Markov's inequality, recalled}
% =====================================================================

The next two sections lean on a single tail bound, proved in
Appendix~\ref{chap:toolkit}: for a nonnegative random variable $X$ and any
$a>0$,
\[
  \Prob(X\ge a)\le\frac{\E[X]}{a}.
\]
In words --- if $\E[X]$ is small, then $X$ is small typically, not merely on
average. This is what lets the alteration method control not just the expected
number of defects but the probability of having too many at once.

% =====================================================================
\section{High girth and high chromatic number}
% =====================================================================

A $k$-clique forces chromatic number at least $k$, a lower bound certified by
purely local information. Is high chromatic number \emph{always} locally
certifiable? The following construction says emphatically no: a graph can be
locally indistinguishable from a tree --- no short cycles at all --- and still
require many colours.

\begin{definition}{Girth}{girth}
The \emph{girth} of a graph is the length of its shortest cycle (infinite if the
graph is a forest).
\end{definition}

\begin{theorem}{Erd\H{o}s 1959}{girth-chromatic}
For every $k$ and $\ell$, there is a graph with girth greater than $\ell$ and
chromatic number greater than $k$.
\end{theorem}

\begin{proof}
Take the Erd\H{o}s--R\'enyi random graph $G\sim G(n,p)$ on $n$ vertices, each
edge present independently with probability $p=(\log n)^2/n$. We control two
features and then delete.

\smallskip\noindent\emph{Short cycles.} Let $X$ count the cycles of length at
most $\ell$. Since $K_n$ has $\binom ni(i-1)!/2\le n^i$ cycles of length $i$,
linearity gives
\[
  \E[X]=\sum_{i=3}^{\ell}\binom ni\frac{(i-1)!}{2}\,p^i
  \le\sum_{i=3}^{\ell}(np)^i
  =\sum_{i=3}^{\ell}(\log n)^{2i}=o(n),
\]
as $\ell$ is constant. By Markov, $\Prob(X\ge n/2)\le \E[X]/(n/2)=o(1)$.

\smallskip\noindent\emph{Large independent sets.} Recall
$\chi(G)\ge |V(G)|/\alpha(G)$, where $\alpha$ is the independence number, so it
suffices to keep $\alpha$ small. With $x=\lceil 3n/\log n\rceil=(3/p)\log n$,
\[
  \Prob\bigl(\alpha(G)\ge x\bigr)\le\binom nx(1-p)^{\binom x2}
  < n^x e^{-p x(x-1)/2}
  = \bigl(n\,e^{-p(x-1)/2}\bigr)^x = o(1),
\]
because $p(x-1)/2\sim\tfrac32\log n$ makes the base $n\cdot n^{-3/2}=n^{-1/2}$.

\smallskip\noindent\emph{Alteration.} For $n$ large, both bad events have
probability below $1/2$, so some $G$ has fewer than $n/2$ short cycles and
$\alpha(G)\le 3n/\log n$. Delete one vertex from each cycle of length
$\le\ell$ to obtain $G'$. Then $|V(G')|\ge n/2$, the girth exceeds $\ell$,
$\alpha(G')\le\alpha(G)\le 3n/\log n$, and
\[
  \chi(G')\ge\frac{|V(G')|}{\alpha(G')}\ge\frac{n/2}{3n/\log n}
  =\frac{\log n}{6}>k
\]
once $n$ is large enough.
\end{proof}

\begin{insight}
This is the alteration method at full power, and the choice $p=(\log n)^2/n$ is
delicate: $p$ must be small enough that short cycles are rare
($np$ only polylogarithmic) yet large enough that independent sets are tiny.
The conclusion is conceptually startling --- high chromatic number is a
\emph{global} property, invisible to any bounded neighbourhood, since a girth
exceeding $\ell$ means every ball of radius $\ell/2$ is a tree. Erd\H{o}s (1962)
sharpened this: there are graphs of chromatic number $>k$ in which \emph{every}
linear-sized subgraph is $3$-colourable, so even seeing a constant fraction of
the graph need not reveal its true colourability.
\end{insight}

% =====================================================================
\section{Random greedy colouring}
% =====================================================================

In Chapter~1 we showed (Theorem~\ref{thm:hyper_lower}) that any $k$-uniform
hypergraph with fewer than $2^{k-1}$ edges is $2$-colourable. We now improve the
threshold by a factor of order $\sqrt{k/\log k}$, via a random greedy colouring.

\begin{theorem}{Radhakrishnan--Srinivasan 2000}{greedy-colouring}
There is a constant $c>0$ such that every $k$-uniform hypergraph with at most
$c\sqrt{\dfrac{k}{\log k}}\;2^{k}$ edges is $2$-colourable.
\end{theorem}

\begin{proof}
Order the vertices by a uniformly random permutation and colour greedily from
left to right: make each vertex blue unless doing so would complete an
all-blue edge, in which case make it red. By construction no edge is all blue,
so the colouring fails only if some edge ends up all red.

\smallskip\noindent\emph{Conflicting pairs.} If a red edge appears, its last
vertex $v$ was forced red, meaning $v$ already completed an all-blue edge $e$
(so $v$ is the \emph{last} vertex of $e$) while also being the \emph{first}
vertex of the red edge $f$. Thus there are two edges $e,f$ sharing the single
vertex $v$, with $v$ last in $e$ and first in $f$; call such an ordered pair
\emph{conflicting}. Crucially, whether $(e,f)$ conflicts depends only on the
random ordering, not on the colours. For two edges meeting in exactly one
vertex (so $2k-1$ vertices in all), conflict requires $v$ to follow the other
$k-1$ vertices of $e$ and precede the other $k-1$ of $f$:
\[
  \Prob\bigl((e,f)\text{ conflicts}\bigr)=\frac{(k-1)!^2}{(2k-1)!}
  \asymp k^{-1/2}2^{-2k}.
\]
Summing over the at most $m^2$ such ordered pairs, the expected number of
conflicting pairs is $\lesssim m^2 k^{-1/2}2^{-2k}$, which is below $1$ once
$m\asymp k^{1/4}2^k$. This already gives a bound (Pluh\'ar 2009); we now sharpen
it.

\smallskip\noindent\emph{The $L\cup M\cup R$ refinement.} Instead of a bare
permutation, drop each vertex at an independent uniform position in $[0,1]$, and
split the interval (with $p$ to be chosen) into
\[
  L=\Bigl[0,\tfrac{1-p}2\Bigr),\quad
  M=\Bigl[\tfrac{1-p}2,\tfrac{1+p}2\Bigr],\quad
  R=\Bigl(\tfrac{1+p}2,1\Bigr].
\]
A fixed edge lands entirely in $L$ with probability $(\tfrac{1-p}2)^k$, and
likewise for $R$, so
\[
  \Prob(\text{some edge lies in }L\text{ or }R)\le 2m\Bigl(\tfrac{1-p}2\Bigr)^k.
\]
Suppose no edge lies wholly in $L$ or $R$. If $(e,f)$ conflicts with shared
vertex $v$, then every vertex of $e$ precedes $v$ and every vertex of $f$
follows it; were $v\in L$ then $e\subseteq L$, and were $v\in R$ then
$f\subseteq R$ --- both excluded. So $v\in M$, and
\[
  \Prob\bigl((e,f)\text{ conflicts, }v\in M\bigr)
  =\int_{(1-p)/2}^{(1+p)/2}x^{k-1}(1-x)^{k-1}\,dx
  \le p\,4^{-k+1},
\]
using $x(1-x)\le\tfrac14$. Union-bounding over $\le m^2$ pairs and combining,
\[
  \Prob(\text{some conflict})
  \le 2m\Bigl(\tfrac{1-p}2\Bigr)^k + m^2 p\,4^{-k+1}
  < 2^{-k+1}m\,e^{-pk} + \bigl(2^{-k+1}m\bigr)^2 p.
\]
Setting $p=\dfrac{1}{k}\log\!\dfrac{2^{k-1}k}{m}$ to minimize the right-hand
side makes it
\[
  \frac{m^2}{4^{k-1}k}\Bigl(1+\log\tfrac{2^{k-1}k}{m}\Bigr),
\]
which is $<1$ when $m=c\,2^k\sqrt{k/\log k}$ for a small constant $c>0$ (taking
$k$ large enough that $p\in[0,1]$; small $k$ are absorbed by shrinking $c$).
With positive probability there is then no conflicting pair, so the greedy
colouring produces no red edge and succeeds.
\end{proof}

\begin{insight}
Where did the improvement come from? The expected number of conflicting pairs
is essentially unchanged --- the $L\cup M\cup R$ device does not make conflicts
rarer on average. What it does is \emph{cluster} the bad events. By first paying
a cheap union bound to forbid any edge from sitting in an extreme block, the
remaining conflicts are confined to a shared common vertex in the thin middle
strip $M$, and the second-order term picks up a factor $p$ (the width of $M$)
that a naive count would miss. Exploiting positive correlation among bad events
that share a localized cause --- rather than bounding them as if independent ---
is the engine of the gain, and squeezing more out of this idea remains an
intriguing open problem.
\end{insight}

% =====================================================================
\section*{Exercises}
\addcontentsline{toc}{section}{Exercises}
% =====================================================================

\begin{enumerate}[itemsep=4pt]
  \item Using the alteration method, prove the Ramsey bound
        $R(4,k)\ge c\,(k/\log k)^2$ for some constant $c>0$.
  \item Prove that every $3$-uniform hypergraph with $n$ vertices and $m\ge n$
        edges contains an independent set (a vertex set spanning no edge) of size
        at least $c\,n^{3/2}/\sqrt m$, for a constant $c>0$.
  \item Prove that every $k$-uniform hypergraph with $n$ vertices and $m$ edges
        has a transversal (a vertex set meeting every edge) of size at most
        $n(\log k)/k+m/k$.
  \item (Zarankiewicz.) Prove that for all $n\ge k\ge2$ there is an $n\times n$
        $\{0,1\}$-matrix with at least $\tfrac12 n^{2-2/(k+1)}$ ones containing no
        all-ones $k\times k$ submatrix.
  \item Fix $k$. Prove there is a constant $c_k>1$ so that for all large $n$
        there exist at least $c_k^{\,n}$ subsets of $[n]$ such that for every $k$
        distinct members $F_1,\dots,F_k$, all $2^k$ intersections
        $\bigcap_i G_i$ are nonempty, where each $G_i$ is $F_i$ or its complement.
  \item (Acute sets.) Prove that for some $c>0$ and every $n$ there are at least
        $c\,(2/\sqrt3)^n$ subsets of $[n]$ with no three distinct members
        $A,B,C$ satisfying $A\cap B\subseteq C\subseteq A\cup B$; deduce a set of
        $\ge c\,(2/\sqrt3)^n$ points in $\mathbb R^n$ spanning only acute angles.
  \item $\bigstar$ Covering complements of sparse graphs by cliques.
        \begin{enumerate}
          \item Prove every $n$-vertex graph with minimum degree $n-d$ is a union
                of $O(d^2\log n)$ cliques.
          \item Prove every bipartite graph with $n$ vertices per side and minimum
                degree $n-d$ ($d\ge1$) is a union of $O(d\log n)$ complete
                bipartite graphs.
        \end{enumerate}
  \item $\bigstar$ Let $G$ have $n$ vertices and minimum degree $\delta\ge2$.
        Prove there is $A\subseteq V$ with $|A|=O(n(\log\delta)/\delta)$ such that
        every vertex outside $A$ has a neighbour in $A$ \emph{and} a neighbour
        outside $A$.
  \item $\bigstar$ Prove every graph $G$ without isolated vertices has an induced
        subgraph on at least $\alpha(G)/2$ vertices in which every vertex has odd
        degree.
\end{enumerate}
